本题不是平均失水曲线拟合，也不是烘房温度优化。待求量是圆柱内部的温度场与干基含水率场。毛细多孔体干燥可用温度与湿度的耦合传递来描述\cite{luikov1975}；本题附录已给出物性与对流系数，因此以题给公式为主模型，不把文献中的全部交叉项搬入。附件提供的是外部环境和几何轨迹，没有药材内部观测，因此模型有效性依靠守恒结构、解析退化核验、网格与时间加密以及物理界限检查，而不能声称已经得到实验标定。

药材长度与直径之比 $L/(2R)=6.25$。两端面积与侧面积之比为 $R/L=0.08$，预热阶段热扩散特征长度小于中截面到端面的距离，因此第一问以中部横截面的径向一维模型作为主线。该近似在数小时乃至数日尺度上会变弱，后文将其作为明确局限而不是已经证明可以忽略的事实。

\subsection{问题一的分析}

第一问时间仅为 $1800\,\mathrm{s}$，热物性取常数，水分扩散系数只依赖于含水率。常物性热方程存在 Bessel--Duhamel 解析核验；变扩散系数的水分问题一般没有同形式闭式解。表面换热与传质均为第三类边界，表面温度和含水率都是待求量，不能直接赋为烘房值。空气含湿量与药材干基含水率即使同以 $\mathrm{kg/kg}$ 表示，质量基准也不一定相同，本文将其作为题给等效传质势使用。

\subsection{问题二的分析}

第二问要求覆盖整个烘干过程的前 $3\,\mathrm{h}$，并改用随含水率变化的密度、比热、导热系数以及同时依赖含水率和温度的扩散系数。预热结束时刻的 $1800\,\mathrm{s}$ 只是第一问的输出期限，并不等于烘房已经进入 $50^\circ\mathrm{C}$ 恒温，因此第二问从 $t=0$ 用新物性重新积分，不拼接第一问数值场。两场通过物性相互耦合，必须联合求解。

\subsection{问题三的分析}

第三问在同一组变物性方程上继续积分，直到药材各处含水率低于 $0.15\,\mathrm{kg/kg}$。``各处''应取计算域内最大值，不能用截面平均或表面值代替。附件环境数据只到 $4\,\mathrm{h}$，之后需要明确的恒温延续假设。后期扩散系数随含水率下降而显著减小，不能用初始扩散率或前 $3\,\mathrm{h}$ 斜率外推总时长。

\subsection{问题四的分析}

第四问同时更换物性并引入半径收缩。附件给出的是外半径轨迹，内部变形不能唯一识别。本文采用均匀比例径向收缩，并引入材料坐标 $\xi=r/R(t)$，使计算域固定在单位区间。干基含水率不是体积浓度，收缩本身不会自动把 $C$ 放大；输出必须按当前物理半径取值，域外单元格不得填零。
